\section{Sparse matrix vector multiplication}
\label{sec:sparse}

SpMV is the family where the v1 report was most wrong, so it is the family that
got rebuilt rather than extended. Two things were broken. The denominator was
contaminated: the cuSPARSE wrapper built its descriptors and sized and allocated
its workspace inside the timed region, so the vendor paid a per call setup my
kernels never paid, and the percentages that came out of that (283, 174 and 127
percent of cuSPARSE) were nonsense in my favour. And the benchmark ran one
synthetic skewed matrix, which cannot separate the kernels that this family exists
to separate.

\subsection{A matrix suite, not a size list}

A sparse kernel's performance is a function of the degree distribution, not of the
dimension, so a size list can only ever confirm the winner that a single matrix
already picked. The family now runs over a suite of real matrices named in
\texttt{experiments/matrix\_manifest.csv}, plus one seeded generator, and every row
of the summary names its matrix rather than a shape. Matrices are fetched ahead of
time by \texttt{scripts/fetch\_matrices.py}; nothing downloads a matrix mid sweep,
because a download inside a measurement is a measurement of the network.

\subsection{The plan, and the structural fix for the timing defect}

\texttt{ckl::SpmvPlan} is built once per matrix and owns everything that does not
change between calls: the cuSPARSE handle, the CSR and dense vector descriptors, a
workspace for each vendor algorithm, the analysis the vendor's second algorithm
needs, the histogram of nonzeros per row, the lane width that histogram implies,
and the SELL-C-sigma and BSR conversions. Running a variant through the plan
enqueues kernel launches and nothing else.

That makes the v1 defect impossible through the public path rather than merely
fixed: there is no entry point in 1.1.0 that builds a descriptor per call. And the
claim is checked rather than asserted. \texttt{bench\_spmv} arms a CUPTI runtime
callback on the allocation and free entry points around its timed lambda and exits
non-zero if the counter moves. A \texttt{-{}-gate-red} switch restores the v1 per
call allocation on purpose so the gate can be watched failing; with the allocation
back in, every variant reads about fifty percent slower, which is the size of the
effect that got into the v1 report.

\subsection{The variants and what each one fixes}

Six kernels, each reachable by name with \texttt{chosen} reporting what ran.

\emph{Naive} is one thread per row, the CSR scalar baseline of Bell and
Garland~\cite{bell2009}: a long row makes its thread run far longer than its warp
neighbours, and the warp retires at the speed of its slowest row. \emph{Warp} is
one 32 lane warp per row with a shuffle reduction, the vector variant of the same
paper: a long row is split 32 ways and the reads coalesce within a row.
\emph{Vector} narrows that to a lane group of 2, 4, 8, 16 or 32 chosen from the
plan's histogram and templated on the width, which fixes the lane waste the warp
variant pays on regular matrices, where a full warp idles 25 lanes and still runs
five reduction steps to combine nothing.

\emph{Merge} is Merrill and Garland's merge based decomposition~\cite{merrill2016spmv}.
The merge of the row end offsets with the nonzero indices has exactly $m + nnz$
steps whatever the degree distribution, so cutting it into equal diagonal segments
gives every thread exactly $\lceil (m + nnz) / \text{threads} \rceil$ units of
work, and a second pass adds in the partials of rows that straddle a segment
boundary. That is the imbalance the warp variant still carries when it hands one
4700 nonzero row to a single warp while its neighbours finish four element rows.

\emph{Sell} is SELL-C-sigma~\cite{kreutzer2014} with C equal to 32, one warp per
slice and one lane per row, rows sorted by length inside a window of sigma rows and
padded to the slice maximum, so the inner loop is a 128 byte column read with no
row pointer indirection at all. It fixes the uncoalesced walk both CSR kernels do
across row boundaries and pays for it in padding, which the summary records as a
padding ratio so the trade is visible rather than implied. \emph{Bsr} is block CSR,
the stretch item: an $r$ by $r$ dense block shares one column index and one span of
the vector across $r^2$ values. It exists for the two matrices with dense blocking
and does nothing for the graph classes, and its baseline is the vendor's block
routine rather than the general one.

\subsection{What is measured}

Table \ref{tab:spmv} is the suite table. It is generated from whatever SpMV rows
the committed summary carries, keyed on the matrix name where the schema has one.
The v1 rows do not name a matrix and their percent of cuSPARSE is withdrawn under
correction A2, so the generator prints the withdrawal rather than the number: the
denominator those percentages were divided by included setup work, and no honest
figure can be recovered from them after the fact. What does survive from v1 is the
kernel to kernel comparison, because both hand written variants were timed
identically: the warp per row kernel is about 1.6 times the thread per row kernel
on a skewed degree matrix. The suite level result, and every percent of cuSPARSE,
is pending the 1.1.0 sweep.

\begin{table}[htbp]
\centering
\begin{tabular}{lllrr}
\toprule
kernel & matrix & nonzeros & GFLOP/s & percent of cuSPARSE \\
\midrule
CSR SpMV naive & synthetic, 16384 rows & pending & 24.7 & withdrawn (A2) \\
CSR SpMV warp & synthetic, 16384 rows & pending & 61.7 & withdrawn (A2) \\
CSR SpMV naive & synthetic, 65536 rows & pending & 48.2 & withdrawn (A2) \\
CSR SpMV warp & synthetic, 65536 rows & pending & 88.8 & withdrawn (A2) \\
CSR SpMV naive & synthetic, 131072 rows & pending & 57.5 & withdrawn (A2) \\
CSR SpMV warp & synthetic, 131072 rows & pending & 93.6 & withdrawn (A2) \\
\bottomrule
\end{tabular}

\caption{CSR SpMV rows in the committed summary. Percent of cuSPARSE is withdrawn
for the v1 rows under correction A2 (setup work inside the vendor's timed region);
the matrix suite result is pending.}
\label{tab:spmv}
\end{table}
